% Options for packages loaded elsewhere
\PassOptionsToPackage{unicode}{hyperref}
\PassOptionsToPackage{hyphens}{url}
\documentclass[
  ignorenonframetext,
]{beamer}
\newif\ifbibliography
\usepackage{pgfpages}
\setbeamertemplate{caption}[numbered]
\setbeamertemplate{caption label separator}{: }
\setbeamercolor{caption name}{fg=normal text.fg}
\beamertemplatenavigationsymbolsempty
% remove section numbering
\setbeamertemplate{part page}{
  \centering
  \begin{beamercolorbox}[sep=16pt,center]{part title}
    \usebeamerfont{part title}\insertpart\par
  \end{beamercolorbox}
}
\setbeamertemplate{section page}{
  \centering
  \begin{beamercolorbox}[sep=12pt,center]{section title}
    \usebeamerfont{section title}\insertsection\par
  \end{beamercolorbox}
}
\setbeamertemplate{subsection page}{
  \centering
  \begin{beamercolorbox}[sep=8pt,center]{subsection title}
    \usebeamerfont{subsection title}\insertsubsection\par
  \end{beamercolorbox}
}
% Prevent slide breaks in the middle of a paragraph
\widowpenalties 1 10000
\raggedbottom
\AtBeginPart{
  \frame{\partpage}
}
\AtBeginSection{
  \ifbibliography
  \else
    \frame{\sectionpage}
  \fi
}
\AtBeginSubsection{
  \frame{\subsectionpage}
}
\usepackage{iftex}
\ifPDFTeX
  \usepackage[T1]{fontenc}
  \usepackage[utf8]{inputenc}
  \usepackage{textcomp} % provide euro and other symbols
\else % if luatex or xetex
  \usepackage{unicode-math} % this also loads fontspec
  \defaultfontfeatures{Scale=MatchLowercase}
  \defaultfontfeatures[\rmfamily]{Ligatures=TeX,Scale=1}
\fi
\usepackage{lmodern}
\usetheme[]{Montpellier}
\ifPDFTeX\else
  % xetex/luatex font selection
\fi
% Use upquote if available, for straight quotes in verbatim environments
\IfFileExists{upquote.sty}{\usepackage{upquote}}{}
\IfFileExists{microtype.sty}{% use microtype if available
  \usepackage[]{microtype}
  \UseMicrotypeSet[protrusion]{basicmath} % disable protrusion for tt fonts
}{}
\makeatletter
\@ifundefined{KOMAClassName}{% if non-KOMA class
  \IfFileExists{parskip.sty}{%
    \usepackage{parskip}
  }{% else
    \setlength{\parindent}{0pt}
    \setlength{\parskip}{6pt plus 2pt minus 1pt}}
}{% if KOMA class
  \KOMAoptions{parskip=half}}
\makeatother
\usepackage{color}
\usepackage{fancyvrb}
\newcommand{\VerbBar}{|}
\newcommand{\VERB}{\Verb[commandchars=\\\{\}]}
\DefineVerbatimEnvironment{Highlighting}{Verbatim}{commandchars=\\\{\}}
% Add ',fontsize=\small' for more characters per line
\usepackage{framed}
\definecolor{shadecolor}{RGB}{248,248,248}
\newenvironment{Shaded}{\begin{snugshade}}{\end{snugshade}}
\newcommand{\AlertTok}[1]{\textcolor[rgb]{0.94,0.16,0.16}{#1}}
\newcommand{\AnnotationTok}[1]{\textcolor[rgb]{0.56,0.35,0.01}{\textbf{\textit{#1}}}}
\newcommand{\AttributeTok}[1]{\textcolor[rgb]{0.13,0.29,0.53}{#1}}
\newcommand{\BaseNTok}[1]{\textcolor[rgb]{0.00,0.00,0.81}{#1}}
\newcommand{\BuiltInTok}[1]{#1}
\newcommand{\CharTok}[1]{\textcolor[rgb]{0.31,0.60,0.02}{#1}}
\newcommand{\CommentTok}[1]{\textcolor[rgb]{0.56,0.35,0.01}{\textit{#1}}}
\newcommand{\CommentVarTok}[1]{\textcolor[rgb]{0.56,0.35,0.01}{\textbf{\textit{#1}}}}
\newcommand{\ConstantTok}[1]{\textcolor[rgb]{0.56,0.35,0.01}{#1}}
\newcommand{\ControlFlowTok}[1]{\textcolor[rgb]{0.13,0.29,0.53}{\textbf{#1}}}
\newcommand{\DataTypeTok}[1]{\textcolor[rgb]{0.13,0.29,0.53}{#1}}
\newcommand{\DecValTok}[1]{\textcolor[rgb]{0.00,0.00,0.81}{#1}}
\newcommand{\DocumentationTok}[1]{\textcolor[rgb]{0.56,0.35,0.01}{\textbf{\textit{#1}}}}
\newcommand{\ErrorTok}[1]{\textcolor[rgb]{0.64,0.00,0.00}{\textbf{#1}}}
\newcommand{\ExtensionTok}[1]{#1}
\newcommand{\FloatTok}[1]{\textcolor[rgb]{0.00,0.00,0.81}{#1}}
\newcommand{\FunctionTok}[1]{\textcolor[rgb]{0.13,0.29,0.53}{\textbf{#1}}}
\newcommand{\ImportTok}[1]{#1}
\newcommand{\InformationTok}[1]{\textcolor[rgb]{0.56,0.35,0.01}{\textbf{\textit{#1}}}}
\newcommand{\KeywordTok}[1]{\textcolor[rgb]{0.13,0.29,0.53}{\textbf{#1}}}
\newcommand{\NormalTok}[1]{#1}
\newcommand{\OperatorTok}[1]{\textcolor[rgb]{0.81,0.36,0.00}{\textbf{#1}}}
\newcommand{\OtherTok}[1]{\textcolor[rgb]{0.56,0.35,0.01}{#1}}
\newcommand{\PreprocessorTok}[1]{\textcolor[rgb]{0.56,0.35,0.01}{\textit{#1}}}
\newcommand{\RegionMarkerTok}[1]{#1}
\newcommand{\SpecialCharTok}[1]{\textcolor[rgb]{0.81,0.36,0.00}{\textbf{#1}}}
\newcommand{\SpecialStringTok}[1]{\textcolor[rgb]{0.31,0.60,0.02}{#1}}
\newcommand{\StringTok}[1]{\textcolor[rgb]{0.31,0.60,0.02}{#1}}
\newcommand{\VariableTok}[1]{\textcolor[rgb]{0.00,0.00,0.00}{#1}}
\newcommand{\VerbatimStringTok}[1]{\textcolor[rgb]{0.31,0.60,0.02}{#1}}
\newcommand{\WarningTok}[1]{\textcolor[rgb]{0.56,0.35,0.01}{\textbf{\textit{#1}}}}
\setlength{\emergencystretch}{3em} % prevent overfull lines
\providecommand{\tightlist}{%
  \setlength{\itemsep}{0pt}\setlength{\parskip}{0pt}}
%\documentclass[unicode,10pt]{beamer}% 'unicode'が必要
\usepackage{luatexja}% 日本語を使う場合は必須
%\usepackage[japanese]{babel}	
%\usefonttheme[onlymath]{serif}
%\usepackage[T1]{fontenc} %これを使うと、ギリシャ文字が出ない
%\usepackage{textcomp}
%\usepackage[scale=1.0]{tgheros} % Sans serif font
%\usepackage[scaled]{beramono} % Monospace font
%\usepackage{luatexja-otf}
%\renewcommand{\kanjifamilydefault}{\gtdefault}
%\usepackage[haranoaji]{luatexja-preset}

% スライド番号の表示 %%%%%%%%%%%%%%%%%%%
\makeatletter
\setbeamertemplate{footline}{%
  \leavevmode%
  \hbox{%
  \begin{beamercolorbox}[wd=.5\paperwidth,ht=2.25ex,dp=1ex,right]{date in head/foot}%
    \insertframenumber{} \hspace*{2ex} % new version without total frames
  \end{beamercolorbox}}%
  \vskip0pt%
}
\makeatother
% スライド番号の表示 %%%%%%%%%%%%%%%%%%%
\usepackage{bookmark}
\IfFileExists{xurl.sty}{\usepackage{xurl}}{} % add URL line breaks if available
\urlstyle{same}
\hypersetup{
  hidelinks,
  pdfcreator={LaTeX via pandoc}}

\title{第8章: 多重回帰分析：不均一分散}
\subtitle{Jeffrey Wooldridge (2016).\\
Introductory Econometrics: A Modern Approach\\
6rd ed.~Cengage Learning.}
\author{}
\date{\vspace{-2.5em}2026-03-05}

\begin{document}
\frame{\titlepage}

\section{準備}\label{ux6e96ux5099}

\begin{frame}[fragile]{必要なパッケージの読み込み}
\phantomsection\label{ux5fc5ux8981ux306aux30d1ux30c3ux30b1ux30fcux30b8ux306eux8aadux307fux8fbcux307f}
\begin{itemize}
\tightlist
\item
  \texttt{wooldridge}パッケージの読み込み
\end{itemize}

\begin{Shaded}
\begin{Highlighting}[]
\FunctionTok{library}\NormalTok{(wooldridge)}
\end{Highlighting}
\end{Shaded}

\begin{itemize}
\tightlist
\item
  \texttt{sandwich}パッケージの読み込み（頑健な標準誤差の計算に使用）
\end{itemize}

\begin{Shaded}
\begin{Highlighting}[]
\FunctionTok{library}\NormalTok{(sandwich)}
\end{Highlighting}
\end{Shaded}

\begin{itemize}
\tightlist
\item
  \texttt{lmtest}パッケージの読み込み（モデルの定式化の検定などに使用）
\end{itemize}

\begin{Shaded}
\begin{Highlighting}[]
\FunctionTok{library}\NormalTok{(lmtest)}
\end{Highlighting}
\end{Shaded}
\end{frame}

\section{8-1
不均一分散の含意}\label{ux4e0dux5747ux4e00ux5206ux6563ux306eux542bux610f}

\begin{frame}{不均一分散 (Heteroskedasticity) とは}
\phantomsection\label{ux4e0dux5747ux4e00ux5206ux6563-heteroskedasticity-ux3068ux306f}
\begin{itemize}
\tightlist
\item
  これまでのガウス＝マルコフ仮定（MLR.5）では、説明変数 \(x\)
  の値に関わらず、誤差項 \(u\) の分散が一定であることを仮定していた。
\item
  \(Var(u|x) = \sigma^2\) （等分散性, Homoskedasticity）
\item
  これが満たされない場合（\(Var(u|x) = \sigma^2(x)\)）、\textbf{不均一分散}と呼ぶ。
\end{itemize}
\end{frame}

\begin{frame}{不均一分散がOLSに与える影響}
\phantomsection\label{ux4e0dux5747ux4e00ux5206ux6563ux304colsux306bux4e0eux3048ux308bux5f71ux97ff}
\begin{enumerate}
\tightlist
\item
  \textbf{不偏性・一致性には影響しない。}
  OLS推定量は依然として不偏かつ一致である。
\item
  \textbf{OLSはもはや最良線形不偏推定量 (BLUE) ではない。}
  より効率的な推定量が存在し得る。
\item
  \textbf{標準誤差の推定が正しくなくなる。} 通常の \(t\) 統計量や \(F\)
  統計量を用いた仮説検定が妥当ではなくなる。
\end{enumerate}
\end{frame}

\section{8-2
不均一分散に頑健な推論}\label{ux4e0dux5747ux4e00ux5206ux6563ux306bux9811ux5065ux306aux63a8ux8ad6}

\begin{frame}{頑健な標準誤差 (Robust Standard Errors)}
\phantomsection\label{ux9811ux5065ux306aux6a19ux6e96ux8aa4ux5dee-robust-standard-errors}
\begin{itemize}
\tightlist
\item
  不均一分散が存在しても、適切な修正を加えることで妥当な推論（\(t\)
  検定など）が可能になる。
\item
  ホワイトの頑健な標準誤差（White standard
  errors）や、不均一分散に頑健な標準誤差（Heteroskedasticity-robust
  standard errors）と呼ばれる。
\item
  大規模標本において、誤差項の分散の形を知らなくても正当化される。
\end{itemize}
\end{frame}

\begin{frame}[fragile]{Rでの実装例 (\texttt{gpa3})}
\phantomsection\label{rux3067ux306eux5b9fux88c5ux4f8b-gpa3}
\footnotesize

\begin{Shaded}
\begin{Highlighting}[]
\FunctionTok{data}\NormalTok{(gpa3)}
\NormalTok{res }\OtherTok{\textless{}{-}} \FunctionTok{lm}\NormalTok{(cumgpa }\SpecialCharTok{\textasciitilde{}}\NormalTok{ sat }\SpecialCharTok{+}\NormalTok{ hsperc }\SpecialCharTok{+}\NormalTok{ tothrs }\SpecialCharTok{+}\NormalTok{ female }\SpecialCharTok{+}\NormalTok{ black }\SpecialCharTok{+}\NormalTok{ white, }
          \AttributeTok{data =}\NormalTok{ gpa3, }\AttributeTok{subset =}\NormalTok{ (spring }\SpecialCharTok{==} \DecValTok{1}\NormalTok{))}

\CommentTok{\# 通常の標準誤差}
\CommentTok{\# summary(res)$coefficients[1:3, 1:2]}

\CommentTok{\# 頑健な標準誤差 (sandwichパッケージを使用)}
\FunctionTok{coeftest}\NormalTok{(res, }\AttributeTok{vcov =} \FunctionTok{vcovHC}\NormalTok{(res, }\AttributeTok{type =} \StringTok{"HC1"}\NormalTok{))[}\DecValTok{1}\SpecialCharTok{:}\DecValTok{3}\NormalTok{, }\DecValTok{1}\SpecialCharTok{:}\DecValTok{2}\NormalTok{]}
\end{Highlighting}
\end{Shaded}

\begin{verbatim}
##                 Estimate   Std. Error
## (Intercept)  1.470064766 0.2206802109
## sat          0.001140728 0.0001915317
## hsperc      -0.008566358 0.0014179275
\end{verbatim}

\normalsize
\end{frame}

\section{8-3
不均一分散の検定}\label{ux4e0dux5747ux4e00ux5206ux6563ux306eux691cux5b9a}

\begin{frame}{ブロイシュ＝ペーガン検定 (Breusch-Pagan Test)}
\phantomsection\label{ux30d6ux30edux30a4ux30b7ux30e5ux30daux30fcux30acux30f3ux691cux5b9a-breusch-pagan-test}
\begin{itemize}
\tightlist
\item
  \(H_0: Var(u|x_1, \dots, x_k) = \sigma^2\) （等分散）
\item
  手順：

  \begin{enumerate}
  \tightlist
  \item
    OLS残差 \(\hat{u}\) を求める。
  \item
    残差の2乗 \(\hat{u}^2\) を独立変数 \(x_j\) で回帰する（補助回帰）。
  \item
    その回帰の \(F\) 統計量または \(LM = n \cdot R^2\)
    を用いて検定する。
  \end{enumerate}
\end{itemize}
\end{frame}

\begin{frame}{ホワイト検定 (White Test)}
\phantomsection\label{ux30dbux30efux30a4ux30c8ux691cux5b9a-white-test}
\begin{itemize}
\tightlist
\item
  BP検定よりも一般的な不均一分散（非線形な形態など）を検出できる。
\item
  手順：残差の2乗 \(\hat{u}^2\)
  を、独立変数、その2乗、および交差項で回帰する。
\item
  変数が多いと自由度が急激に減少するため、予測値 \(\hat{y}\) とその2乗
  \(\hat{y}^2\) を用いる簡略版がよく使われる。
\end{itemize}
\end{frame}

\begin{frame}[fragile]{検定の実装例}
\phantomsection\label{ux691cux5b9aux306eux5b9fux88c5ux4f8b}
\footnotesize

\begin{Shaded}
\begin{Highlighting}[]
\CommentTok{\# BP検定 (lmtestパッケージ)}
\FunctionTok{bptest}\NormalTok{(res)}
\end{Highlighting}
\end{Shaded}

\begin{verbatim}
## 
##  studentized Breusch-Pagan test
## 
## data:  res
## BP = 44.557, df = 6, p-value = 5.732e-08
\end{verbatim}

\begin{Shaded}
\begin{Highlighting}[]
\CommentTok{\# ホワイト検定（簡略版）}
\CommentTok{\# 手動計算: 残差2乗を予測値とその2乗で回帰}
\NormalTok{u2 }\OtherTok{\textless{}{-}} \FunctionTok{resid}\NormalTok{(res)}\SpecialCharTok{\^{}}\DecValTok{2}
\NormalTok{yhat }\OtherTok{\textless{}{-}} \FunctionTok{fitted}\NormalTok{(res)}
\NormalTok{white\_reg }\OtherTok{\textless{}{-}} \FunctionTok{lm}\NormalTok{(u2 }\SpecialCharTok{\textasciitilde{}}\NormalTok{ yhat }\SpecialCharTok{+} \FunctionTok{I}\NormalTok{(yhat}\SpecialCharTok{\^{}}\DecValTok{2}\NormalTok{))}
\CommentTok{\# LM統計量 (n * R2) と p値}
\NormalTok{n }\OtherTok{\textless{}{-}} \FunctionTok{length}\NormalTok{(u2)}
\NormalTok{lm\_stat }\OtherTok{\textless{}{-}}\NormalTok{ n }\SpecialCharTok{*} \FunctionTok{summary}\NormalTok{(white\_reg)}\SpecialCharTok{$}\NormalTok{r.squared}
\NormalTok{p\_val }\OtherTok{\textless{}{-}} \DecValTok{1} \SpecialCharTok{{-}} \FunctionTok{pchisq}\NormalTok{(lm\_stat, }\AttributeTok{df =} \DecValTok{2}\NormalTok{)}
\FunctionTok{data.frame}\NormalTok{(}\AttributeTok{LM =}\NormalTok{ lm\_stat, }\AttributeTok{p\_value =}\NormalTok{ p\_val)}
\end{Highlighting}
\end{Shaded}

\begin{verbatim}
##         LM   p_value
## 1 1.264452 0.5314075
\end{verbatim}

\normalsize
\end{frame}

\section{8-4 加重最小二乗法
(WLS)}\label{ux52a0ux91cdux6700ux5c0fux4e8cux4e57ux6cd5-wls}

\begin{frame}{加重最小二乗法 (Weighted Least Squares)}
\phantomsection\label{ux52a0ux91cdux6700ux5c0fux4e8cux4e57ux6cd5-weighted-least-squares}
\begin{itemize}
\tightlist
\item
  不均一分散の構造 \(Var(u|x) = \sigma^2 h(x)\)
  が既知である場合、OLSよりも効率的な推定が可能。
\item
  各観測値を \(1/\sqrt{h(x)}\)
  倍して重み付けし、変換されたモデルにOLSを適用する。
\item
  これを\textbf{一般化最小二乗法 (GLS)} の一種と呼ぶ。
\end{itemize}
\end{frame}

\begin{frame}{実行可能な一般化最小二乗法 (FGLS)}
\phantomsection\label{ux5b9fux884cux53efux80fdux306aux4e00ux822cux5316ux6700ux5c0fux4e8cux4e57ux6cd5-fgls}
\begin{itemize}
\tightlist
\item
  関数 \(h(x)\) の形が未知の場合、データから \(h(x)\)
  を推定してWLSを行う。
\item
  一般的な手順：

  \begin{enumerate}
  \tightlist
  \item
    \(\log(\hat{u}^2)\) を独立変数で回帰し、予測値 \(\hat{g}\) を得る。
  \item
    重み \(\hat{h} = \exp(\hat{g})\) を計算する。
  \item
    重み \(1/\hat{h}\) を用いてWLSを実行する。
  \end{enumerate}
\end{itemize}
\end{frame}

\begin{frame}{まとめ}
\phantomsection\label{ux307eux3068ux3081}
\begin{itemize}
\tightlist
\item
  不均一分散下では、OLS推定量は不偏だが効率的ではなく、通常の検定が使えない。
\item
  \textbf{解決策1:} 頑健な標準誤差を用いる（最も一般的）。
\item
  \textbf{解決策2:}
  不均一分散の検定を行い、必要に応じてWLS/FGLSで効率性を高める。
\item
  実証分析では、まずは頑健な標準誤差で結果を確認するのが定石である。
\end{itemize}
\end{frame}

\end{document}
